\documentclass[12pt,a4paper]{article}

\usepackage[utf8]{inputenc}
\usepackage[T2A]{fontenc}
\usepackage[russian,english]{babel}
\usepackage{amsmath,amssymb,bm}
\usepackage{graphicx}
\usepackage{float}
\usepackage{booktabs}
\usepackage{geometry}
\usepackage[expansion=false]{microtype}
\usepackage{hyperref}
\usepackage{xcolor}

\geometry{margin=2.2cm}
\setlength{\emergencystretch}{2em}

% Figures are placed at the top of the page, including float-only pages.
\makeatletter
\setlength{\@fptop}{0pt}
\setlength{\@fpsep}{12pt}
\setlength{\@fpbot}{0pt plus 1fil}
\makeatother
\hypersetup{colorlinks=true,linkcolor=black,urlcolor=blue}

\graphicspath{{figures/}}

\newcommand{\reportfigure}[3]{%
\begin{figure}[H]
    \centering
    \includegraphics[width=0.94\textwidth]{#1}%
    \caption{#2}
    \label{#3}
\end{figure}
}

\title{Дисперсия экситон--фотонных поляритонов\\ в планарной DBR-микрополости}
\author{Коростелев Никита}
\date{}

\begin{document}
\maketitle
\clearpage

\section{Цель работы}

Цель работы --- рассчитать дисперсию экситон--фотонных поляритонов в планарной полупроводниковой микрополости и понять, как она возникает из связи света в полости с экситоном в квантовой яме. Аналитические модели используются для понимания происхождения основных зависимостей, а затем сравниваются с расчетом всей многослойной структуры методом матриц переноса.

DBR (distributed Bragg reflector) --- многослойное брэгговское зеркало. QW (quantum well) --- квантовая яма.

Для наклонного распространения отдельно рассматриваются TE- и TM-поляризации. В TE электрическое поле перпендикулярно плоскости падения, а в TM этой плоскости перпендикулярно магнитное поле.

Основная физическая схема состоит из фотонной моды микрополости и экситона в квантовой яме. При слабой связи их энергии почти не влияют друг на друга. При сильной связи возникают новые смешанные состояния --- нижний и верхний поляритоны.

\clearpage
\section{Структура микрополости и параметры}

Рассматривается структура

\[
\mathrm{air}\,|\,(HL)^{23}\,|\,C_1\,|\,W\,|\,C_2\,|\,(LH)^{26}\,|\,\mathrm{air},
\]

где $H$ --- слой GaAs с большим показателем преломления, $L$ --- слой AlAs с меньшим показателем, $C_1$ и $C_2$ --- части GaAs-слоя полости между зеркалами, а $W$ --- квантовая яма In$_{0.08}$Ga$_{0.92}$As. Свет падает слева через верхнее DBR-зеркало. С обеих сторон полости к ее GaAs-слою примыкает слой $L$.

\reportfigure{01_refractive_index_profile.png}{Профиль показателя преломления $n(z)$ многослойной микрополости. Рисунок строится непосредственно по заданной геометрии и показателям преломления слоев, то есть не является результатом решения дисперсионной задачи. Красной линией отмечен центр QW.}{fig:index-profile}

На рис.~\ref{fig:index-profile} видны два конечных DBR-зеркала и более широкий слой полости между ними. Толщина QW мала по сравнению с полной длиной структуры, поэтому на этом масштабе она прежде всего отмечает свое положение внутри полости.

В качестве экспериментального ориентира используется структура из работы P. Cilibrizzi et al., \textit{Polariton condensation in a strain-compensated planar microcavity with InGaAs quantum wells}, Applied Physics Letters 105, 191118 (2014): $2\lambda$-полость, 23 и 26 периодов DBR, квантовые ямы толщиной $6$ нм и суммарное расщепление Раби порядка $8$ мэВ для шести сильно связанных QW.

Для рассматриваемой структуры используются постоянные показатели преломления

\[
n_H=3.60,\qquad n_L=2.95,\qquad n_C=3.60,\qquad n_W=3.65.
\]

Для расчета фиксируется энергия экситонного резонанса при нулевом поперечном импульсе,

\[
E_X(0)=1.450\ \mathrm{eV}.
\]

Здесь $E_X(0)$ --- энергия экситона при $k_\parallel=0$. Значение $1.450$ эВ относится к той же области энергий, около $1.45$ эВ, где Cilibrizzi et al. наблюдали поляритонные ветви. Оно задает центр рассматриваемого спектра и не является результатом нашего расчета.

Соответствующая длина волны в вакууме

\[
\lambda_0=\frac{hc}{E_X(0)}\approx0.855\ \mu\mathrm m.
\]

Начальные толщины четвертьволновых слоев DBR равны

\[
d_H=\frac{\lambda_0}{4n_H}\approx59.4\ \mathrm{nm},
\qquad
 d_L=\frac{\lambda_0}{4n_L}\approx72.5\ \mathrm{nm}.
\]

Для $2\lambda$-полости начальная оценка полной толщины слоя полости имеет вид

\[
L_C^{(0)}=\frac{2\lambda_0}{n_C}\approx475\ \mathrm{nm}.
\]

Конечное значение $L_C$ определяется по резонансу всей DBR-структуры. Причина в том, что отражение от DBR дает дополнительный фазовый сдвиг и немного меняет условие резонанса.

Квантовая яма имеет толщину

\[
d_{QW}=6\ \mathrm{nm}
\]

и заменяет равный по толщине участок GaAs-слоя полости. Поэтому

\[
L_C=d_{C_1}+d_{QW}+d_{C_2}.
\]

\clearpage
\section{Дисперсия невзаимодействующих мод}

\subsection{Фотонная мода микрополости}

В планарной полости продольная компонента волнового вектора фиксируется условием резонанса, тогда как $k_\parallel$ остается непрерывным. Поэтому энергия фотона имеет вид

\[
E_C(k_\parallel)
=
\frac{\hbar c}{n_C}
\sqrt{k_z^2+k_\parallel^2}.
\]

Если использовать энергию моды при нормальном распространении, эту формулу удобно записать как

\[
E_C(k_\parallel)
=
\sqrt{
E_C(0)^2+
\left(\frac{\hbar c}{n_C}k_\parallel\right)^2
}.
\]

При малом $k_\parallel$ получается параболическая дисперсия

\[
E_C(k_\parallel)
\approx
E_C(0)+\frac{\hbar^2k_\parallel^2}{2m_{ph}},
\]

где

\[
m_{ph}=\frac{E_C(0)n_C^2}{c^2}.
\]

Параболическая форма возникает не из-за ненулевой массы свободного фотона, а из-за фиксации $k_z$ полостью: при изменении только поперечного импульса квадратный корень раскладывается в ряд по $k_\parallel^2$.

\subsection{Экситон в квантовой яме}

Квантовая яма ограничивает движение электрона и дырки по $z$, но центр масс экситона остается свободным в плоскости слоя. Поэтому его дисперсия имеет вид

\[
E_X(k_\parallel)
=
E_X(0)+\frac{\hbar^2k_\parallel^2}{2M_X},
\]

где $M_X$ --- масса центра масс экситона. В расчете используется

\[
M_X=0.5m_0.
\]

Так как $M_X\gg m_{ph}$, экситонная ветвь значительно более плоская, чем фотонная. Для аналитических графиков используется тот же диапазон $k_\parallel=0\ldots4\ \mu\mathrm m^{-1}$, что и в дальнейшем численном расчете.

Эти две невзаимодействующие дисперсии далее не выносятся в отдельный рисунок: они показаны вместе с поляритонными ветвями на следующем графике. Это позволяет сразу увидеть, как исходное пересечение фотонной и экситонной мод превращается в антипересечение после включения связи.


\clearpage
\section{Аналитическая модель поляритона}

При фиксированном поперечном волновом векторе $k_\parallel$ рассматриваются две степени свободы: фотон в микрополости и экситон в квантовой яме. В книге Kavokin et al., \textit{Microcavities}, \S\S5.8.1--5.8.2 эта система записывается как две связанные квантовые моды:

\[
\hat H
=
\sum_{k_\parallel}
E_C(k_\parallel)\hat B_{k_\parallel}^{\dagger}\hat B_{k_\parallel}
+
\sum_{k_\parallel}
E_X(k_\parallel)\hat X_{k_\parallel}^{\dagger}\hat X_{k_\parallel}
+
\sum_{k_\parallel}
g(k_\parallel)
\left(
\hat B_{k_\parallel}^{\dagger}\hat X_{k_\parallel}
+
\hat X_{k_\parallel}^{\dagger}\hat B_{k_\parallel}
\right).
\]

Здесь $\hat B_{k_\parallel}$ убирает один фотон из моды полости, а $\hat B_{k_\parallel}^{\dagger}$ создает его. Аналогично, $\hat X_{k_\parallel}$ и $\hat X_{k_\parallel}^{\dagger}$ убирают и создают экситон. Величины $E_C(k_\parallel)$ и $E_X(k_\parallel)$ --- энергии фотона и экситона без связи между ними, а $g$ --- энергия их взаимодействия. Один и тот же $k_\parallel$ у фотона и экситона сохраняется из-за однородности структуры в плоскости слоев.

В рассматриваемом диапазоне связь считается постоянной:

\[
g(k_\parallel)\approx g,
\qquad
\Omega_R=2g,
\]

где $\Omega_R$ --- расщепление Раби: расстояние по энергии между верхним и нижним поляритонами в точке, где исходные фотонная и экситонная моды имели бы одинаковую энергию.

Для каждого фиксированного $k_\parallel$ гамильтониан сводится к матрице

\[
H(k_\parallel)=
\begin{pmatrix}
E_C(k_\parallel) & g\\
g & E_X(k_\parallel)
\end{pmatrix}
=
\begin{pmatrix}
E_C(k_\parallel) & \Omega_R/2\\
\Omega_R/2 & E_X(k_\parallel)
\end{pmatrix}.
\]

Расстройка невзаимодействующих мод при произвольном $k_\parallel$ определяется как

\[
\Delta(k_\parallel)
=
E_C(k_\parallel)-E_X(k_\parallel).
\]

Отдельно удобно обозначить расстройку при нормальном распространении

\[
\delta
=
E_C(0)-E_X(0).
\]

Если обе невзаимодействующие дисперсии записать в параболическом приближении, то

\[
\Delta(k_\parallel)
=
\delta
+
\frac{\hbar^2k_\parallel^2}{2}
\left(
\frac{1}{m_{ph}}-\frac{1}{M_X}
\right),
\]

где $m_{ph}$ --- эффективная масса фотона полости, а $M_X$ --- масса центра масс экситона в плоскости QW. Поскольку масса экситона $M_X$ намного больше эффективной фотонной массы $m_{ph}$, экситонная ветвь меняется медленно. Поэтому изменение расстройки с $k_\parallel$ почти целиком задается фотонной ветвью.

Собственные энергии находятся из условия

\[
\det\!\left[H(k_\parallel)-E I\right]=0,
\]

то есть

\[
\left[E_C(k_\parallel)-E\right]
\left[E_X(k_\parallel)-E\right]
-g^2=0.
\]

Два решения этого уравнения образуют верхнюю и нижнюю поляритонные ветви:

\[
E_{UP,LP}(k_\parallel)
=
\frac{E_C(k_\parallel)+E_X(k_\parallel)}{2}
\pm
\frac12
\sqrt{\Delta(k_\parallel)^2+\Omega_R^2}.
\]

Индекс $UP$ относится к знаку ``$+$'' и обозначает верхнюю поляритонную ветвь, индекс $LP$ --- к знаку ``$-$'' и обозначает нижнюю поляритонную ветвь. При точном резонансе, $\Delta=0$,

\[
E_{UP}-E_{LP}=\Omega_R.
\]

Поэтому пересечение невзаимодействующих мод заменяется антипересечением.

\reportfigure{02_analytical_polariton_dispersion.png}{Аналитическая дисперсия фотона, экситона и поляритонов при нулевой расстройке. Все четыре кривые получены аналитически: $E_C(k_\parallel)$ и $E_X(k_\parallel)$ задают невзаимодействующие моды, а $E_{LP}(k_\parallel)$ и $E_{UP}(k_\parallel)$ --- собственные значения матрицы $2\times2$. По вертикальной оси энергия показана относительно $E_X(0)$.}{fig:analytical-dispersion}

На рис.~\ref{fig:analytical-dispersion} видно, что связь заменяет пересечение исходных мод антипересечением. При увеличении $k_\parallel$ нижняя ветвь приближается к почти плоскому экситону, а верхняя все сильнее следует фотонной дисперсии.

Поляритон является линейной комбинацией фотонного и экситонного состояний. В обозначениях \S5.8.1 книги преобразование Хопфилда имеет вид

\[
\hat a_{LP,k}
=
X_k\hat X_k-C_k\hat B_k,
\qquad
\hat a_{UP,k}
=
C_k\hat X_k+X_k\hat B_k,
\]

Здесь $\hat a_{LP,k}$ и $\hat a_{UP,k}$ относятся к нижнему и верхнему поляритонам. Величины $C_k$ и $X_k$ называются коэффициентами Хопфилда. Они нормированы условием

\[
|C_k|^2+|X_k|^2=1.
\]

Величина $|C|^2$ показывает фотонную долю поляритона, а $|X|^2$ --- экситонную. Для нижней ветви

\[
|C_{LP}(k_\parallel)|^2
=
\frac12
\left(
1-
\frac{\Delta(k_\parallel)}
{\sqrt{\Delta(k_\parallel)^2+\Omega_R^2}}
\right),
\]

\[
|X_{LP}(k_\parallel)|^2
=
\frac12
\left(
1+
\frac{\Delta(k_\parallel)}
{\sqrt{\Delta(k_\parallel)^2+\Omega_R^2}}
\right).
\]

Для верхней ветви доли меняются местами:

\[
|C_{UP}(k_\parallel)|^2=|X_{LP}(k_\parallel)|^2,
\qquad
|X_{UP}(k_\parallel)|^2=|C_{LP}(k_\parallel)|^2.
\]

При $\Delta=0$ обе поляритонные моды имеют равные фотонную и экситонную доли,

\[
|C|^2=|X|^2=\frac12.
\]

При большой расстройке одна ветвь становится в основном фотонной, а другая --- в основном экситонной.

\reportfigure{03_hopfield_coefficients.png}{Фотонная и экситонная доли нижнего и верхнего поляритонов. График рассчитан аналитически из собственных векторов той же модели $2\times2$; величины $|C|^2$ и $|X|^2$ показывают соответственно фотонную и экситонную доли состояния.}{fig:hopfield}

Рис.~\ref{fig:hopfield} дает простой способ читать дальнейшие результаты: с ростом $k_\parallel$ нижний поляритон становится преимущественно экситонным, а верхний --- преимущественно фотонным. Эта смена состава далее проявится и в отражении, и в ширинах линий.

\clearpage
\section{Конечная DBR-микрополость}

Аналитическая парабола не учитывает конечное число слоев зеркал, фазовый сдвиг при отражении и различие TE/TM при ненулевом $k_\parallel$. Поэтому далее используется метод матриц переноса (TMM, transfer-matrix method). Его смысл простой: для каждого слоя отдельно учитываются распространение волны и отражение на границах, после чего матрицы всех слоев перемножаются. Общая постановка приведена в Kavokin et al., \textit{Microcavities}, §§2.5--2.7.

Для слоя с номером $j$, показателем преломления $n_j$ и толщиной $d_j$ продольная компонента волнового вектора равна

\[
q_j(E,k_\parallel)
=
\sqrt{
\left(\frac{n_jE}{\hbar c}\right)^2-k_\parallel^2
}.
\]

При прохождении слоя фаза волны меняется. Для слоя толщины $d_j$ это изменение фазы равно

\[
\varphi_j=q_jd_j.
\]

Здесь $\varphi_j$ показывает, на сколько радиан меняется фаза при прохождении слоя $j$.

Чтобы одной формулой записывать условия на границах слоев для TE и TM, вводятся вспомогательные величины $Y_j$:

\[
Y_j^{TE}=n_j\cos\phi_j,
\qquad
Y_j^{TM}=\frac{\cos\phi_j}{n_j},
\]

где $\phi_j$ --- угол распространения внутри слоя $j$, а

\[
\cos\phi_j=
\frac{q_j}{n_jE/(\hbar c)}.
\]

Здесь $Y_j$ не является новой физической характеристикой материала: это удобная комбинация $n_j$ и угла $\phi_j$, которая позволяет одинаково записать граничные условия для TE и TM. Эти зависимости соответствуют §2.5 книги \textit{Microcavities}.

В базисе прямой и обратной волн

\[
\mathbf v(z)=
\begin{pmatrix}
A^+(z)\\A^-(z)
\end{pmatrix}
\]

матрица распространения через однородный слой имеет вид

\[
P_j=
\begin{pmatrix}
e^{iq_jd_j} & 0\\
0 & e^{-iq_jd_j}
\end{pmatrix}.
\]

Матрица границы $i\rightarrow j$ равна

\[
I_{i\rightarrow j}
=
\frac{1}{2Y_j}
\begin{pmatrix}
Y_i+Y_j & Y_j-Y_i\\
Y_j-Y_i & Y_i+Y_j
\end{pmatrix}.
\]

Полная матрица $M$ связывает амплитуды слева и справа:

\[
\mathbf v_R=M\mathbf v_L.
\]

При падении света слева

\[
\mathbf v_L=
\begin{pmatrix}
1\\r
\end{pmatrix},
\qquad
\mathbf v_R=
\begin{pmatrix}
t\\0
\end{pmatrix}.
\]

Отсюда

\[
r=-\frac{M_{21}}{M_{22}},
\qquad
R=|r|^2.
\]

Сначала влияние экситона на свет выключается, но сам слой InGaAs остается в полости. Так определяется фотонный резонанс той же геометрии без экситон--фотонной связи. Толщина полости $L_C$ подбирается так, чтобы при $k_\parallel=0$

\[
E_C^{bg}(0)=E_X(0)
\]

для нулевой расстройки. Это исключает дополнительный сдвиг фотонной моды из-за отличия $n_W$ от $n_C$.

Для дисперсионных расчетов используется диапазон $k_\parallel=0\ldots4\ \mu\mathrm m^{-1}$ с шагом около $0.0167\ \mu\mathrm m^{-1}$. Исходный спектр отражения рассчитывается по энергии с шагом $0.0005$ мэВ. Чтобы узкие резонансы были хорошо видны на карте, при построении рисунка показывается минимум отражения внутри энергетического интервала шириной $0.050$ мэВ.

Обозначим эту отображаемую величину как

\[
R_{\mathrm{bin}}(E',k_\parallel)
=
\min_{E\in\mathrm{bin}(E')}R_{TE}(E,k_\parallel),
\]

где $E'$ --- центр соответствующего энергетического интервала. Такое бинирование используется только для визуализации и не меняет сам расчет структуры.

\reportfigure{04_bare_cavity_reflectivity_map.png}{Численно рассчитанный минимум TE-отражения внутри энергетических интервалов шириной $0.050$ мэВ для микрополости с выключенным экситонным резонансом. Узкая резонансная линия соответствует фотонной моде конечной DBR-структуры.}{fig:bare-reflectivity}

На рис.~\ref{fig:bare-reflectivity} видна одна фотонная ветвь. Ее плавный рост с $k_\parallel$ соответствует обычной дисперсии моды планарной полости, но здесь эта ветвь уже получена не из аналитической формулы, а из полного расчета слоев.

Для количественного сравнения фотонная энергия в конечной DBR-структуре определяется не по идеальной формуле, а из собственной моды открытой полости. При принятой матричной записи она удовлетворяет условию

\[
M_{22}(\mathcal E_C^{bg},k_\parallel)=0.
\]

Численная фотонная дисперсия задается действительной частью полюса $\operatorname{Re}\mathcal E_C^{bg}(k_\parallel)$. Сравнение с идеальной полостью показывает влияние конечных DBR и их фазового отклика на форму ветви.

\reportfigure{05_bare_photon_dispersion_comparison.png}{Фотонная дисперсия: аналитическая модель идеальной плоской полости и численный расчет конечной DBR-структуры. Пунктирная кривая получена аналитически из $E_C(k_\parallel)$; сплошная --- из действительной части комплексного полюса $\mathcal E_C^{bg}$, найденного методом матриц переноса из условия $M_{22}=0$. Энергия показана относительно $E_X(0)$.}{fig:bare-dispersion-comparison}

На рис.~\ref{fig:bare-dispersion-comparison} обе модели почти совпадают при малом $k_\parallel$, но затем численная ветвь становится заметно более крутой. Это показывает влияние конечных DBR: фаза отражения и проникновение поля в зеркала меняются с энергией и углом, чего нет в формуле идеальной полости.

\clearpage
\section{Квантовая яма в микрополости}

Квантовая яма $W$ имеет конечную толщину $d_{QW}=6$ нм и показатель преломления $n_W$. Сам слой InGaAs влияет на распространение света как обычный диэлектрик. Дополнительно внутри него есть экситонный переход: около энергии $E_X$ квантовая яма особенно сильно взаимодействует со светом. Для этого резонансного отклика используется модель одиночной QW из \S\S4.4.2 и 4.4.4 книги \textit{Microcavities}.

Энергия экситона при заданном $k_\parallel$ снова равна

\[
E_X(k_\parallel)
=
E_X(0)
+
\frac{\hbar^2k_\parallel^2}{2M_X},
\]

где $E_X(0)$ --- энергия экситонного резонанса при $k_\parallel=0$, а $M_X$ --- масса центра масс экситона в плоскости квантовой ямы.

Для слоя InGaAs продольная компонента волнового вектора имеет вид

\[
q_W(E,k_\parallel)
=
\sqrt{
\left(\frac{n_WE}{\hbar c}\right)^2-k_\parallel^2
}.
\]

Угол распространения внутри QW обозначим $\phi_W$. Тогда

\[
\cos\phi_W
=
\frac{q_W}{n_WE/(\hbar c)}.
\]

Амплитуда отражения одиночной QW имеет резонансный вид; см. \textit{Microcavities}, \S4.4.2. В энергетических единицах используем

\[
r_X(E,k_\parallel)
=
\frac{i\Gamma_0^{pol}}
{E_X(k_\parallel)-E-i\left(\Gamma_0^{pol}+\gamma_X\right)}.
\]

Здесь $r_X$ --- амплитуда отраженной волны, создаваемой экситонным резонансом. Величина $\Gamma_0^{pol}$ задает радиационное уширение для выбранной поляризации, а $\gamma_X$ --- дополнительное однородное уширение экситона из-за процессов в материале. Обе величины записаны как HWHM (half width at half maximum), то есть \textit{полуширина линии на половине ее максимума}. Поэтому суммарная полуширина резонанса равна

\[
\gamma_{QW}^{pol}
=
\Gamma_0^{pol}+\gamma_X,
\]

а полная ширина линии на половине максимума, FWHM (full width at half maximum), равна

\[
\Gamma_{QW}^{pol}=2\gamma_{QW}^{pol}.
\]

Далее для интуитивного обсуждения достаточно помнить: увеличение $\Gamma_0^{pol}$ или $\gamma_X$ делает экситонный резонанс шире.

Амплитуда прохождения связана с амплитудой отражения соотношением

\[
t_X=1+r_X,
\]

где $t_X$ --- амплитудный коэффициент прохождения резонансной QW.

При наклонном распространении радиационный параметр зависит от поляризации. Согласно \S4.4.4 книги,

\[
\Gamma_0^{TE}
=
\frac{\Gamma_0}{\cos\phi_W},
\qquad
\Gamma_0^{TM}
=
\Gamma_0\cos\phi_W,
\]

где $\Gamma_0$ --- радиационная полуширина линии при нормальном распространении. При $k_\parallel=0$ имеем $\phi_W=0$, поэтому

\[
\Gamma_0^{TE}=\Gamma_0^{TM}=\Gamma_0.
\]

Резонансный отклик экситона включается в тот же метод матриц переноса через матрицу

\[
T_X
=
\frac{1}{t_X}
\begin{pmatrix}
t_X^2-r_X^2 & r_X\\
-r_X & 1
\end{pmatrix}.
\]

Здесь $T_X$ описывает только экситонный резонанс. Распространение через конечную толщину InGaAs учитывается матрицей

\[
P_W(\ell)
=
\begin{pmatrix}
e^{iq_W\ell} & 0\\
0 & e^{-iq_W\ell}
\end{pmatrix},
\]

где $\ell$ --- пройденная толщина QW. Для экситонного листа, помещенного в центре слоя толщины $d_{QW}$, полный блок квантовой ямы равен

\[
M_{QW}
=
I_{W\rightarrow C}
P_W(d_{QW}/2)
T_X
P_W(d_{QW}/2)
I_{C\rightarrow W}.
\]

Матрицы $I_{C\rightarrow W}$ и $I_{W\rightarrow C}$ описывают границы GaAs--InGaAs и InGaAs--GaAs. Если экситонный резонанс выключен, то

\[
T_X=I,
\]

но слой InGaAs с показателем $n_W$ и толщиной $d_{QW}$ остается в структуре. Такая структура задает фотонный резонанс $E_C^{bg}$. Индекс $bg$ означает background: учитывается обычный показатель преломления QW, но ее экситонный резонанс выключен.

Квантовая яма заменяет равный по толщине участок GaAs-слоя полости. Поэтому полная толщина полости сохраняется:

\[
L_C=d_{C_1}+d_{QW}+d_{C_2}.
\]

Если $z_{QW}$ --- координата центра квантовой ямы, отсчитанная от левой границы полости, то

\[
d_{C_1}=z_{QW}-\frac{d_{QW}}{2},
\]

\[
d_{C_2}=L_C-z_{QW}-\frac{d_{QW}}{2}.
\]

Таким образом, изменение положения QW меняет $d_{C_1}$ и $d_{C_2}$, но не увеличивает полную толщину $L_C$.

Для качественного сравнения тот же физический смысл можно передать более простой моделью двух связанных резонансов. С учетом потерь ее удобно записать как

\[
\left(E_X-E-i\gamma_X\right)
\left(E_C-E-i\frac{\Gamma_C}{2}\right)
=g^2,
\]

где $E_C$ --- энергия фотонного резонанса без связи с экситоном, $\Gamma_C$ --- полная ширина фотонной линии на половине максимума, $E_X$ --- энергия экситона, а $g$ --- энергия экситон--фотонной связи. В пределе малых потерь

\[
\Omega_R=2g.
\]

Эта связь показывает физический смысл полной модели: экситонный и фотонный резонансы снова образуют две связанные ветви, как в матрице $2\times2$. Отличие состоит в том, что метод матриц переноса одновременно учитывает реальные зеркала, положение QW внутри стоячей волны и поляризацию.

Параметр $\Gamma_0$ определяется по величине связи одной QW со светом. Для численного ориентира используется тот же эксперимент Cilibrizzi et al. с $2\lambda$ GaAs-микрополостью и In$_{0.08}$Ga$_{0.92}$As квантовыми ямами. В нем три пары квантовых ям расположены около пучностей поля. В каждой двойке находятся две QW, поэтому число сильно связанных ям равно $3\times2=6$. Еще две QW в том же образце расположены около минимумов поля и служат главным образом для сбора носителей, поэтому в это число не входят. Для шести сильно связанных QW наблюдалось суммарное расщепление Раби около $8$ мэВ. Если отдельные QW находятся в примерно одинаковом поле, коллективная связь растет приблизительно как $\sqrt{N}$, поэтому ориентир для одной QW равен

\[
\Omega_R^{(1)}
\approx
\frac{8}{\sqrt{6}}\ \mathrm{meV}
\approx3.27\ \mathrm{meV}.
\]

Здесь $N=6$ получается именно из трех пар QW: $N=3\times2$. Величина $\Omega_R^{(1)}$ --- оценка расщепления, которое дала бы одна такая QW в сопоставимом поле. Параметр $\Gamma_0$ затем настраивается так, чтобы при нулевой расстройке, то есть при $E_C(0)=E_X(0)$, и при $k_\parallel=0$ выполнялось

\[
\operatorname{Re}\mathcal E_{UP}(0)
-
\operatorname{Re}\mathcal E_{LP}(0)
=
3.27\ \mathrm{meV},
\]

где $\mathcal E_{LP}$ и $\mathcal E_{UP}$ --- энергии нижней и верхней поляритонных мод с учетом их конечного времени жизни; смысл комплексной записи подробно вводится ниже.

\clearpage
\section{Поляритонная дисперсия методом матриц переноса}

После включения экситонного отклика $T_X$ одна фотонная линия в спектре отражения разделяется на две поляритонные ветви. Нижняя ветвь обозначается $LP$, верхняя --- $UP$. Их форма зависит и от дисперсий экситона и фотона, и от распространения света через конечные DBR-зеркала.

Для этой карты используется та же сетка по $k_\parallel$ и исходная энергетическая сетка, что и выше. Для отображения применяется более узкий энергетический интервал $0.025$ мэВ, поскольку теперь нужно различить две близкие поляритонные особенности.

\reportfigure{06_polariton_reflectivity_map_te.png}{Численно рассчитанный минимум TE-отражения внутри энергетических интервалов шириной $0.025$ мэВ после включения экситонного резонанса QW. Две резонансные особенности соответствуют нижней и верхней поляритонным ветвям.}{fig:polariton-reflectivity}

На рис.~\ref{fig:polariton-reflectivity} одна фотонная линия из рис.~\ref{fig:bare-reflectivity} превращается в две ветви. Нижняя ветвь с ростом $k_\parallel$ становится слабее в отражении, что согласуется с рис.~\ref{fig:hopfield}: ее фотонная доля уменьшается.

Далее численный результат сравнивается с моделью $2\times2$. В этом сравнении важно не вводить второй источник различий через приближенную фотонную параболу. Поэтому фотонная энергия в матрице $2\times2$ берется из той же конечной многослойной структуры с выключенным экситонным резонансом:

\[
E_C(k_\parallel)=\operatorname{Re}\mathcal{E}_C^{bg}(k_\parallel).
\]

Экситонная энергия по-прежнему задается аналитически как

\[
E_X(k_\parallel)=E_X(0)+\frac{\hbar^2k_\parallel^2}{2M_X}.
\]

После этого собственные значения матрицы $2\times2$ сравниваются с минимумами отражения полной структуры. Таким образом, на этом рисунке проверяется прежде всего сама модель связи фотона с экситоном: фотонная мода в обоих подходах относится к одной и той же конечной DBR-микрополости.

\reportfigure{07_polariton_tmm_vs_2x2.png}{Поляритонная дисперсия: модель $2\times2$ и численный метод матриц переноса. Линии получены аналитической диагонализацией матрицы $2\times2$, но в качестве $E_C(k_\parallel)$ использована численно найденная фотонная мода фоновой DBR-структуры $\operatorname{Re}\mathcal E_C^{bg}$. Кружки --- положения минимумов $R_{TE}$ в полном расчете с экситонным резонансом.}{fig:tmm-vs-2x2}

На рис.~\ref{fig:tmm-vs-2x2} точки почти лежат на линиях модели $2\times2$. Это означает, что в выбранном диапазоне простая модель двух связанных мод хорошо описывает само экситон--фотонное взаимодействие, если фотонную дисперсию брать из реальной конечной полости.

\clearpage
\section{Поле в полости и положение QW}

Положение QW влияет на силу связи с фотоном: связь больше там, где поле стоячей волны сильнее. Для полости без экситон--фотонной связи при $k_\parallel=0$ восстанавливается стоячая волна. Если на левой границе слоя $j$ известны амплитуды $A_j^+$ и $A_j^-$, то внутри слоя

\[
E_j(\xi)
=
A_j^+e^{iq_j\xi}
+
A_j^-e^{-iq_j\xi}.
\]

Для сравнения используется нормированная интенсивность

\[
I(z)=
\frac{|E(z)|^2}
{\max_{\mathrm{cavity}}|E(z)|^2}.
\]

Профиль поля восстанавливается численно с шагом $\Delta z=0.5$ нм. Одиночная QW помещается в пучность, ближайшую к центру полости, а затем смещается относительно нее.

Чтобы отделить влияние положения QW от изменения расстройки, при каждом смещении сначала рассматривается та же структура с выключенным экситонным резонансом и фотонный резонанс снова настраивается на прежнее значение $E_C^{bg}(0)$. После включения экситонного отклика из двух поляритонных полюсов определяется реальное расстояние между модами

\[
\Delta E_{\mathrm{pole}}(\Delta z_{QW})
=
\operatorname{Re}
\left[
\mathcal E_{UP}(0)-\mathcal E_{LP}(0)
\right],
\]

где $\Delta z_{QW}$ --- смещение центра квантовой ямы относительно выбранной пучности поля.

Для одной QW есть простое ожидание: энергия связи пропорциональна амплитуде электрического поля в месте расположения ямы,

\[
g(z_{QW})\propto |E(z_{QW})|.
\]

Поскольку $I(z)\propto|E(z)|^2$, оценку величины $2g$ можно записать как

\[
2g_{\mathrm{field}}(\Delta z)=\Omega_1\sqrt{\frac{I(z_0+\Delta z)}{I(z_0)}}.
\]

где $\Omega_1$ --- расщепление одной QW в выбранной пучности, а $z_0$ --- положение этой пучности. Эта формула не подставляется в численный расчет: она строится отдельно и сравнивается с $\Delta E_{\mathrm{pole}}$.

\reportfigure{08_qw_position_coupling.png}{Положение QW и сила экситон--фотонной связи. Верхняя панель --- численно восстановленная нормированная интенсивность $I(z)$ фотонной моды; точки показывают пучности, красная линия --- выбранное положение QW. Нижняя панель сравнивает оценку $2g_{\mathrm{field}}\propto\sqrt{I}$ с реальным расстоянием между действительными частями поляритонных полюсов, полученным из полного расчета методом матриц переноса.}{fig:qw-position}

На рис.~\ref{fig:qw-position} обе зависимости почти совпадают в большей части диапазона. Это прямо показывает, что сила связи определяется локальной амплитудой поля. Только совсем близко к узлам поля реальное расстояние между полюсами может уменьшаться быстрее простой оценки, потому что при очень слабой связи заметнее становится влияние потерь.

\clearpage
\section{Влияние числа пар DBR}

Конечная прозрачность DBR приводит к утечке фотона из полости. При увеличении числа пар отражение возрастает, резонанс становится уже, а время жизни фотонной моды увеличивается.

Для комплексной энергии фотонной моды, где $E_C$ обозначает положение резонанса,

\[
\mathcal E_C=E_C-i\frac{\Gamma_C}{2}
\]

ширина линии равна

\[
\Gamma_C=-2\operatorname{Im}\mathcal E_C.
\]

Соответствующая добротность полости равна

\[
Q_C=\frac{E_C}{\Gamma_C}.
\]

Поэтому уменьшение ширины линии одновременно означает рост добротности и времени жизни фотонной моды.

Для сравнения рассматриваются, например,

\[
(N_t,N_b)=(8,8),\ (12,12),\ (16,16),\ (23,26).
\]

Для каждого числа пар полость заново настраивается на ту же энергию при $k_\parallel=0$, поэтому изменение $\Gamma_C$ связано именно с зеркалами.

В этой серии учитывается только утечка через конечные идеальные DBR. Остальные потери не моделируются.

Поэтому результат описывает прозрачность зеркал, а не полную экспериментальную ширину линии.

\reportfigure{09_dbr_linewidth_sweep.png}{Ширина фотонной моды, связанная с утечкой через конечные DBR, в зависимости от числа пар зеркал. Все точки получены численно из комплексного фотонного полюса: $\Gamma_C^{\mathrm{mirror}}=-2\operatorname{Im}\mathcal E_C$. Логарифмическая шкала показывает быстрое уменьшение потерь при увеличении числа пар.}{fig:dbr-linewidth}

Рис.~\ref{fig:dbr-linewidth} показывает ожидаемую тенденцию: чем больше пар в DBR, тем слабее фотон выходит через зеркала и тем уже резонанс. Для структуры $(23,26)$ расчетная ширина очень мала именно потому, что модель здесь описывает идеализированную зеркальную утечку, а не все потери экспериментального образца.

\clearpage
\section{TE--TM расщепление}

При нормальном распространении направления TE и TM эквивалентны. Поэтому

\[
E^{TE}(0)=E^{TM}(0).
\]

При $k_\parallel\neq0$ фазовые свойства DBR и отклик QW становятся поляризационно-зависимыми. В \textit{Microcavities}, §§2.7 и 4.4.4, показано, что дисперсии TE и TM могут иметь различную кривизну.

Для поляритонных ветвей определим

\[
\Delta_{TE-TM}^{LP}(k_\parallel)
=
E_{LP}^{TE}(k_\parallel)-E_{LP}^{TM}(k_\parallel),
\]

\[
\Delta_{TE-TM}^{UP}(k_\parallel)
=
E_{UP}^{TE}(k_\parallel)-E_{UP}^{TM}(k_\parallel).
\]

Обе величины получаются из одного и того же численного расчета. На верхней панели удобно показать сами четыре ветви, а на нижней --- их малую разность, которая на общей энергетической шкале почти не видна.

Комплексные полюса TE и TM ищутся на той же сетке $k_\parallel=0\ldots4\ \mu\mathrm m^{-1}$, что и остальные дисперсионные ветви. По энергии отдельная сетка для полюса не нужна: комплексный корень уточняется непосредственно для каждого значения $k_\parallel$.

\reportfigure{10_te_tm_polariton_splitting.png}{TE--TM расщепление поляритонных мод. Верхняя панель показывает численно найденные ветви $E_{LP}^{TE}$, $E_{UP}^{TE}$, $E_{LP}^{TM}$ и $E_{UP}^{TM}$; нижняя --- их малые разности $\Delta_{TE-TM}^{LP}$ и $\Delta_{TE-TM}^{UP}$. Все величины получены из комплексных полюсов той же многослойной структуры.}{fig:te-tm}

На общей шкале верхней панели рис.~\ref{fig:te-tm} TE- и TM-ветви почти совпадают. Нижняя панель показывает, что расщепление все же есть и растет до десятков $\mu$эВ. Для UP знак разности меняется: это означает лишь смену того, какая из двух поляризаций имеет немного большую энергию.

\clearpage
\section{Открытая микрополость: энергия и затухание мод}

DBR-зеркала отражают свет очень хорошо, но не идеально. Поэтому фотон и поляритон могут постепенно выходить из микрополости. Резонанс имеет не только центральную энергию, но и конечную ширину: чем быстрее энергия уходит из полости, тем шире спектральная линия и тем меньше время жизни моды.

Эти две величины удобно объединить в комплексную энергию

\[
\mathcal E(k_\parallel)
=
E(k_\parallel)-i\frac{\Gamma(k_\parallel)}{2}.
\]

Здесь $E=\operatorname{Re}\mathcal E$ --- положение резонанса по энергии, а $\Gamma=-2\operatorname{Im}\mathcal E$ --- полная ширина линии на половине максимума. Соответствующее характерное время жизни равно

\[
\tau(k_\parallel)=\frac{\hbar}{\Gamma(k_\parallel)}.
\]

Чтобы найти такую моду, мысленно убирают падающую извне волну и ищут поле, которое может существовать в полости само по себе и только уходить наружу через зеркала. Для принятой матричной записи это условие имеет вид

\[
M_{22}(\mathcal E,k_\parallel)=0.
\]

Комплексные полюса рассчитываются на той же сетке $k_\parallel$ с шагом около $0.0167\ \mu\mathrm m^{-1}$. В отличие от карты отражения, по энергии здесь не перебирается отдельная сетка: для каждого $k_\parallel$ комплексный корень уравнения уточняется напрямую. Такое описание открытых мод обсуждается в Kavokin et al., \textit{Microcavities}, \S2.6.

Для нижнего и верхнего поляритонов получаются две комплексные энергии $\mathcal E_{LP}$ и $\mathcal E_{UP}$. Их действительные части дают уже знакомые ветви $E_{LP}$ и $E_{UP}$, а мнимые части дают ширины

\[
\Gamma_{LP,UP}(k_\parallel)
=-2\operatorname{Im}\mathcal E_{LP,UP}(k_\parallel).
\]

Таким образом, вдоль одной и той же дисперсионной ветви можно одновременно видеть, как меняется энергия поляритона и насколько быстро он затухает. Отдельный рисунок с действительными частями комплексных энергий здесь не нужен: те же дисперсионные ветви уже показаны выше. В этом разделе интерес представляет прежде всего их мнимая часть, то есть ширина линии.

Для сравнения с простой моделью $2\times2$ потери можно добавить прямо в ее матрицу:

\[
H=
\begin{pmatrix}
E_C-i\Gamma_C/2 & \Omega_R/2\\
\Omega_R/2 & E_X-i\Gamma_X^{nr}/2
\end{pmatrix}.
\]

Здесь $\Gamma_C$ --- ширина фотонного резонанса, а $\Gamma_X^{nr}=2\gamma_X$ --- нерадиационная ширина экситонного перехода. Эта матрица помогает понять простую физику: поляритон, который в данной области больше похож на фотон, сильнее наследует фотонные потери; более экситонная мода сильнее наследует экситонное уширение.

В пределе слабых потерь это можно приблизительно выразить через коэффициенты Хопфилда:

\[
\Gamma_{LP}
\approx
|C_{LP}|^2\Gamma_C
+
|X_{LP}|^2\Gamma_X^{nr},
\]

\[
\Gamma_{UP}
\approx
|C_{UP}|^2\Gamma_C
+
|X_{UP}|^2\Gamma_X^{nr}.
\]

Эти формулы нужны для физической интерпретации. Сами кривые ширин ниже не строятся из этой оценки: они получаются непосредственно из мнимых частей комплексных полюсов полной многослойной структуры.

\reportfigure{11_polariton_linewidths.png}{Ширины линий нижнего и верхнего поляритонов. Обе кривые получены численно из комплексных полюсов полной структуры, $\Gamma_{LP,UP}=-2\operatorname{Im}\mathcal E_{LP,UP}$. Оценка через коэффициенты Хопфилда используется только для физической интерпретации.}{fig:polariton-linewidths}

Связь с рис.~\ref{fig:hopfield} здесь особенно наглядна. При $k_\parallel=0$ обе ветви имеют близкий состав и близкие ширины. Затем LP становится экситонным и его ширина приближается к экситонной, а UP становится фотонным и его ширина уменьшается к узкой ширине фотонной моды из рис.~\ref{fig:dbr-linewidth}.

Минимумы отражения и комплексные полюса относятся к одной и той же открытой микрополости, но извлекаются из расчета разными способами. Поэтому отдельно проверяется, насколько энергия спектрального провала совпадает с действительной частью соответствующего полюса.

\reportfigure{12_reflectivity_dips_vs_poles.png}{Сравнение минимумов отражения с действительными частями комплексных полюсов. Кружки получены из минимумов численно рассчитанного $R_{TE}(E,k_\parallel)$, а линии --- из решения $M_{22}(\mathcal E,k_\parallel)=0$ для той же многослойной структуры.}{fig:dips-vs-poles}

Рис.~\ref{fig:dips-vs-poles} показывает, что провалы отражения хорошо отслеживают действительные части поляритонных полюсов. Поэтому точки и линии на рис.~\ref{fig:tmm-vs-2x2} также хорошо согласуются.
\clearpage
\section{Коллективная связь нескольких квантовых ям}

До этого рассматривалась одна активная QW. В экспериментальном образце Cilibrizzi et al. использовались три пары квантовых ям, расположенные около пучностей стоячего поля. Поэтому со светом сильно взаимодействовали шесть QW. Здесь этот случай используется только как дополнительная проверка коллективной связи; точная эпитаксиальная структура каждой двойки QW не воспроизводится.

Сначала рассмотрим простой физический смысл. Если одна QW связана с фотонной модой с энергией взаимодействия $g_1$, то для нескольких одинаковых QW, взаимодействующих с одной и той же фотонной модой, эффективная связь равна

\[
g_N=
\sqrt{g_1^2+g_2^2+\ldots+g_N^2}.
\]

Если все ямы находятся в близком по величине поле, то $g_i\approx g_1$, и поэтому

\[
g_N\approx g_1\sqrt{N_{QW}},
\]

а расщепление Раби должно расти как

\[
\Omega_R(N_{QW})\approx\Omega_R(1)\sqrt{N_{QW}}.
\]

Это аналитическое ожидание. Оно показывает, почему несколько QW усиливают связь: все они взаимодействуют с одной фотонной модой совместно.

В численной модели рассматриваются $N_{QW}=1,2,4,6$ эквивалентных активных QW. Они размещаются в областях сильного поля, найденных по профилю на рис.~\ref{fig:qw-position}. Точная структура барьера внутри каждой экспериментальной двойки QW здесь не задается: две близкие ямы около одной пучности учитываются как два одинаковых экситонных вклада в одной области поля. Общая толщина полости при этом сохраняется.

Для каждого $N_{QW}$ сначала выключается экситонный отклик и фотонный резонанс снова настраивается на нулевую расстройку. После этого экситонный отклик включается и при $k_\parallel=0$ находятся два комплексных поляритонных полюса. Численное расщепление определяется как

\[
\Omega_R^{\mathrm{num}}(N_{QW})
=
\operatorname{Re}
\left[
\mathcal E_{UP}(0)-\mathcal E_{LP}(0)
\right].
\]

Здесь отдельная энергетическая сетка не используется: два полюса находятся напрямую для каждого из четырех значений $N_{QW}$.

\reportfigure{13_multiple_qw_scaling.png}{Коллективное усиление связи при нескольких квантовых ямах. Кружки --- расщепление $\Omega_R^{\mathrm{num}}$, найденное численно из полюсов многослойной структуры для $N_{QW}=1,2,4,6$. Линия --- аналитическая зависимость $\Omega_R(1)\sqrt{N_{QW}}$.}{fig:multi-qw}

На рис.~\ref{fig:multi-qw} численный расчет следует зависимости $\sqrt{N_{QW}}$. Это показывает, что модель воспроизводит ожидаемое коллективное усиление связи.

Важно, что значение для одной QW было ранее выбрано из экспериментального ориентира $8/\sqrt{6}\approx3.27$ мэВ. Поэтому точка около $8$ мэВ при $N_{QW}=6$ не является независимым предсказанием экспериментального результата. Здесь проверяется именно закон изменения связи при увеличении числа активных QW.

\clearpage
\section{Заключение}

Аналитическая картина последовательно сравнивается с моделью всей многослойной микрополости. Модель $2\times2$ объясняет антипересечение и состав LP/UP (рис.~\ref{fig:analytical-dispersion} и \ref{fig:hopfield}). Метод матриц переноса показывает те же эффекты в конечной DBR-структуре.

Сравнение на рис.~\ref{fig:bare-dispersion-comparison} показывает, что конечные зеркала заметно меняют фотонную дисперсию при больших $k_\parallel$. При этом после подстановки этой численной фотонной ветви простая модель $2\times2$ почти совпадает с полным расчетом поляритонов (рис.~\ref{fig:tmm-vs-2x2}). Значит, основная физика сильной связи действительно хорошо описывается двумя связанными модами.

Рис.~\ref{fig:qw-position} показывает прямую связь между амплитудой стоячего поля и силой взаимодействия с QW. Рис.~\ref{fig:polariton-linewidths} связывает изменение ширин линий со сменой фотонной и экситонной долей поляритона. Для нескольких QW численный расчет воспроизводит ожидаемый закон коллективного усиления $\Omega_R\propto\sqrt{N_{QW}}$ (рис.~\ref{fig:multi-qw}).

Таким образом, аналитические формулы дают физическую интерпретацию, а численная модель учитывает конечные DBR, положение QW, поляризацию и открытость микрополости.

\end{document}
